% ════════════════════════════════════════════════════════════════
\subsection{Termin 2}
% ════════════════════════════════════════════════════════════════

\begin{zadanie}[Zadanie 1 --- scalanie list (15 p.)]
	\textbf{(i) (3p.)} Niech $A,B$ będą listami długości $k$ i~$l$ posortowanymi
	rosnąco. Podaj algorytm działający w~czasie $\BO(k+l)$, który scala je w~jedną
	posortowaną rosnąco listę długości $k+l$.\\
	\textbf{(ii) (4p.)} Mamy $m$ list długości $n$ każda, wszystkie posortowane
	rosnąco; scalamy je w~jedną długości $mn$, używając wielokrotnie algorytmu
	z~(i): scalamy dwie pierwsze listy, wynik z~trzecią, ten z~czwartą, itd.
	Oszacuj złożoność jako funkcję $m$ i~$n$.\\
	\textbf{(iii) (8p.)} Podaj asymptotycznie szybszy algorytm typu ,,dziel
	i~zwyciężaj'' rozwiązujący problem z~(ii). Oszacuj jego złożoność.
\end{zadanie}

\begin{zadanie}[Zadanie 2 --- brak słów kodowych długości~1 (15 p.)]
	Załóżmy, że w~tekście, który chcemy zakodować za pomocą kodu Huffmana, każdy
	znak występuje mniej niż $\frac{24}{3}$ razy. Udowodnij, że w~kodzie Huffmana
	dla takiego tekstu żaden znak nie będzie zakodowany za pomocą tylko jednego
	bitu.
\end{zadanie}

\begin{zadanie}[Zadanie 3 --- pakowanie do skrzyń (Next-Fit) jest 2-aproksymacyjne (15 p.)]
	Rozważmy następujący algorytm dla problemu pakowania do skrzyń. Wejściem jest
	nieposortowana lista rozmiarów pakowanych obiektów. Algorytm pakuje obiekty
	w~kolejności z~listy, próbując umieścić kolejny obiekt tylko w~ostatnio
	rozpoczętej skrzyni; jeśli to niemożliwe, rozpoczyna nową skrzynię. Udowodnij,
	że jest to algorytm $2$-aproksymacyjny dla problemu pakowania do skrzyń.
\end{zadanie}

\begin{zadanie}[Zadanie 4 --- najdłuższy wspólny podciąg kolejnych wyrazów (15 p.)]
	Dla zadanych ciągów $(x_1,x_2,\dots,x_n)$ i~$(y_1,y_2,\dots,y_m)$ chcemy
	znaleźć długość ich najdłuższego wspólnego podciągu \emph{kolejnych} wyrazów,
	czyli największe $k$, dla którego istnieją indeksy $i,j$ takie, że
	$(x_i,x_{i+1},\dots,x_{i+k-1})=(y_j,y_{j+1},\dots,y_{j+k-1})$. Podaj algorytm
	programowania dynamicznego o~złożoności $\BO(mn)$ rozwiązujący ten problem.
\end{zadanie}

